%!TEX root = main.tex

In this paper, motivated by developing a formally verified compiler for the {\firrtl} language, we have proposed a denotational formal semantics of {\firrtl} programs. Specifically, we define the semantics of a {\firrtl} program as the set of low-level RTL circuits that conform to it. 
%The definition of the semantics is technically challenging since we need to tackle the abstraction-level gap between low-level circuits and high-level programs. 
With the formal semantics, we then formalized the three key lowering transformation passes of {\hifirrtl} programs, namely, ``InferWidths'', ``ExpandConnects'', and ``ExpandWhens''. Furthermore, we prove the functional correctness of these three passes, that is, the semantics is preserved by these passes. %For the ``InferWidths'' pass, we actually proved more than the preservation of the semantics. We show that the widths inferred by the ``InferWidths'' pass are the minimum ones among the feasible widths. 
Finally, to validate our approach, we extracted an OCaml compiler out of the coq implementation of the three passes and compared the OCaml compiler with {\firtool}, the official {\firrtl} compiler, on 113 programs. The experiment results confirms the effectiveness and efficiency of our compiler. 

%show the pros of the OCaml compiler on the correctness with respect to {\firtool}: While the OCaml compiler produces 112 programs that are semantically equivalent with those produced by {\firtool}, {\firtool} fails to compile one of them. Moreover, the passes implemented in the OCaml compiler turn out to be more efficient  than those implemented in {\firtool} on average.  

This work paves the way towards a fully fledged and formally verified compiler for {\firrtl} in the future, where more compilation passes e.g. various optimizations can be implemented and their correctness can be verified. 